\exposetitle{XXIV}{Automorphisms of reductive groups}
\label{XXIV}
% original PDF p. 1
\setcounter{footnote}{-1}
\nde{Version of 13/10/2024.}The first part of this exposé (nos. 1 to 5) is a direct consequence of the existence, for a reductive group, of ``sufficiently many outer automorphisms'', a result which is a consequence of the weakest form of the isomorphism theorem for pinned groups. The second part (nos. 6 and 7) presents two applications of the more precise results of the preceding exposé; in particular, no. 7 uses the theorem on generators and relations in its explicit form. Finally, in an appendix (no. 8) we have given results on ``Galois'' cohomology used in the text.

Let us specify our cohomological notation: if $S$ is a scheme and $G$ an $S$-group scheme, we shall denote by $\mathrm H^1(S,G)$ the first cohomology set of $S$ with coefficients in $G$, computed for the (fpqc) topology; it is also the set of isomorphism classes of (fpqc) principal homogeneous sheaves under $G$. We shall denote by $\mathrm H^1_{\mathrm{\acute et}}(S,G)$ the corresponding set for the étale topology; it is therefore the part of $\mathrm H^1(S,G)$ consisting of the classes of homogeneous sheaves under $G$ which are quasi-isotrivial (= locally trivial for the étale topology). We shall denote by $\operatorname{Fib}(S,G)$ the part of $\mathrm H^1(S,G)$ consisting of the classes of representable sheaves (principal homogeneous bundles). We therefore have the inclusions
\[
\mathrm H^1_{\mathrm{\acute et}}(S,G)\subset\mathrm H^1(S,G),\qquad
\operatorname{Fib}(S,G)\subset\mathrm H^1(S,G).
\]
If every principal homogeneous sheaf under $G$ is representable (for example if $G$ is quasi-affine over $S$, cf. SGA~1, VIII 7.9), we therefore have $\operatorname{Fib}(S,G)=\mathrm H^1(S,G)$.

If $S'\to S$ is a covering morphism for the (fpqc) topology, we denote by $\mathrm H^1(S'/S,G)$ the kernel of the canonical map $\mathrm H^1(S,G)\to\mathrm H^1(S',G_{S'})$. One knows that $\mathrm H^1(S'/S,G)$ can be computed simplicially (TDTE~I, §~A.4), which implies that when $S'\to S$ is covering for the étale topology, $\mathrm H^1(S'/S,G)$ is also the kernel of $\mathrm H^1_{\mathrm{\acute et}}(S,G)\to\mathrm H^1_{\mathrm{\acute et}}(S',G_{S'})$.

% original PDF p. 2
Finally, following Exp.~VIII, 4.5, one calls ``theorem 90'' the following assertion: ``every principal homogeneous sheaf under $\mathbf G_{\mathrm m,S}$ is representable and locally trivial'', an assertion equivalent to ``$\mathrm H^1(S,\mathbf G_{\mathrm m,S})=\Pic(S)$'', or again to ``$\mathrm H^1(S,\mathbf G_{\mathrm m,S})=0$ for $S$ local (or, more generally, semi-local)''.

\sgasection{1}{Scheme of automorphisms of a reductive group}
\numpara{1.0}\label{XXIV.1.0}
It is appropriate first to specify certain definitions from the preceding exposé. Let $S$ be a nonempty scheme, $G$ a reductive $S$-group, and $\mathscr R=(M,M^*,R,R^*,\Delta)$ a reduced pinned root datum (Exp.~XXIII 1.5). One calls a \emph{pinning of $G$ of type $\mathscr R$}, or an \emph{$\mathscr R$-pinning of $G$}, the datum:
\begin{enumeratei}
\item of an isomorphism of $\mathrm D_S(M)$ onto a maximal torus $T$ of $G$ (or, equivalently, of a monomorphism $\mathrm D_S(M)\to G$ whose image is a maximal torus $T$ of $G$), identifying $R$ with a root system of $G$ relative to $T$ (Exp.~XIX, 3.6) and $R^*$ with the set of corresponding coroots,
\item for each $\alpha\in\Delta$, of an $X_\alpha\in\Gamma(S,\mathfrak g^\alpha)^\times$.
\end{enumeratei}
For $G$ to possess an $\mathscr R$-pinning, it is necessary and sufficient that it be splittable and of type $\mathscr R$ (Exp.~XXII, 2.7).

If $u:G\to G'$ is an isomorphism of reductive $S$-groups, to every $\mathscr R$-pinning $\mathscr E$ of $G$ there corresponds, by ``transport of structure'', an $\mathscr R$-pinning $u(\mathscr E)$ of $G'$. If $v:\mathscr R'\to\mathscr R$ is an isomorphism of pinned root data, to every $\mathscr R'$-pinning $\mathscr E$ of $G'$ there corresponds by transport of structure an $\mathscr R$-pinning $v(\mathscr E)$ of $G$.
% typo?: the source says G' and G in this transport-of-structure sentence; kept as printed.

Let us call a \emph{pinned group} a triple $(G,\mathscr R,\mathscr E)$ where $G$ is a reductive $S$-group, $\mathscr R$ a reduced pinned root datum, and $\mathscr E$ a pinning of $G$ of type $\mathscr R$. One calls an \emph{isomorphism} of the pinned group $(G,\mathscr R,\mathscr E)$ onto the pinned group $(G',\mathscr R',\mathscr E')$ a pair $(u,v)$ where $u$ is an isomorphism $u:G\to G'$ and $v$ an isomorphism of pinned root data $v:\mathscr R'\to\mathscr R$, such that $u(\mathscr E)=v(\mathscr E')$.\nde{Thus $\Lie(u)(X_{v(\alpha)})=X'_\alpha$, for every $\alpha\in R'$.}

N.~B. \emph{If $S$ is nonempty}, $v$ is uniquely determined by $u$, and by abuse of language one will also say that $u$ is an isomorphism of the pinned groups. In particular, if $(G,\mathscr R,\mathscr E)$ is a pinned group, an automorphism of $(G,\mathscr R,\mathscr E)$ is therefore an automorphism $u$ of $G$ such that there exists an automorphism $v$ of $\mathscr R$ such that $u(\mathscr E)=v(\mathscr E)$; it is therefore an automorphism of $G$, normalizing $T$, inducing on $T$ an automorphism of the form $\mathrm D_S(h)$, where $h$ is an automorphism of $M$, and permuting the elements $X_\alpha$, $\alpha\in\Delta$, among themselves. (As one readily sees, the preceding conditions in fact characterize the automorphisms of $(G,\mathscr R,\mathscr E)$.)

One has an evident contravariant functor
\[
\Phi:\quad (G,\mathscr R,\mathscr E)\mto\mathscr R,\qquad (u,v)\mto v
\]
and the principal result of the preceding exposé (Exp.~XXIII, 4.1) shows us that it is a fully faithful functor (we shall moreover see in the following exposé that it is an equivalence of categories). It follows in particular that the group of automorphisms of $(G,\mathscr R,\mathscr E)$ is canonically isomorphic to the group of automorphisms of the pinned root datum $\mathscr R$ (cf. Exp.~XXIII, 5.5).

% original PDF p. 3
\numpara{1.1}\label{XXIV.1.1}
Let $S$ be a scheme; equip $\Sch/S$ with the (fpqc) topology and consider the $S$-sheaf of groups $\underline{\Aut}_{S\text{-gr.}}(G)$, where $G$ is an $S$-group scheme. One has an exact sequence of $S$-sheaves of groups
\[
\xymatrix{1\ar[r]&\uCentr(G)\ar[r]&G\ar[r]^{\mathrm{int}}&\underline{\Aut}_{S\text{-gr.}}(G)}
\]
which defines a monomorphism
\[
j:G/\uCentr(G)\to\underline{\Aut}_{S\text{-gr.}}(G).
\]
The image sheaf of $j$ is the sheaf of inner automorphisms of $G$; for an automorphism $u$ of $G$ to be inner, it is necessary and sufficient that there exist a covering family $\{S_i\to S\}$ and, for each $i$, a $g_i\in G(S_i)$ such that $\mathrm{int}(g_i)=u_{S_i}$. In this case, if $v$ is another automorphism of $G$, one sees at once that $\mathrm{int}(v)u=vuv^{-1}$ is the inner automorphism defined by the family $g'_i=v(g_i)$. It follows that the image of $j$ is normal in $\underline{\Aut}_{S\text{-gr.}}(G)$. The quotient sheaf of groups, denoted by $\underline{\operatorname{Autext}}(G)$, is the sheaf of outer automorphisms of $G$. One therefore has an exact sequence
\[
1\to G/\uCentr(G)\to\underline{\Aut}_{S\text{-gr.}}(G)
\to\underline{\operatorname{Autext}}(G)\to1.
\]
The preceding definitions are all compatible with base change. They are of course valid in any site.

\numpara{1.2}\label{XXIV.1.2}
Let $S$ be a scheme and $(G,\mathscr R,\mathscr E)$ a pinned reductive $S$-group. Let $E$ be the (abstract) group of automorphisms of the pinned root datum $\mathscr R$, i.e. the group of automorphisms of $\mathscr R$ normalizing $\Delta$. By Exp.~XXIII, 5.5, one has a canonical monomorphism
\[
E\to\underline{\Aut}_{S\text{-gr.}}(G)
\]
which associates to $h\in E$ the unique automorphism $u$ of the pinned group $G$ such that $\Phi(u)=h$. This monomorphism canonically defines a monomorphism of sheaves
\begin{equation}\tag{$*$}
a:E_S\to\underline{\Aut}_{S\text{-gr.}}(G).
\end{equation}
For an automorphism $u$ of $G$ to be a section of the image sheaf of $a$, it is necessary and sufficient that the following conditions hold:
\begin{enumeratei}
\item $u$ normalizes $T$. One then knows that $u$ permutes the roots of $G$ relative to $T$. If $\alpha\in R$, then $u(\alpha):t\mto\alpha(u^{-1}(t))$ is therefore locally on $S$ of the form $t\mto\beta(t)$, with $\beta\in R$. The second condition can then be written as follows:
\item If $\alpha\in\Delta$ and if $U$ is an open subset of $S$ such that $u(\alpha)_U\in R$, then $u(\alpha)_U\in\Delta$ and
\[
\Lie(u_U)(X_\alpha)_U=(X_{u(\alpha)})_U.
\]
\end{enumeratei}
It follows at once from the definitions that the sections of $a(E_S)$ normalize the subgroups of $G$ defined by the pinning: $T$, $B$, $B^-$, $U$, $U^-$. These definitions being given, one has:

% original PDF p. 4
\begin{theorem}{1.3}\label{XXIV.1.3}
Let $S$ be a scheme and $G$ a reductive $S$-group. Consider the canonical exact sequence of $S$-sheaves of groups\nde{Recall that $\ad(G)=G/\uCentr(G)$ denotes the adjoint group of $G$, cf. XXII 4.3.6.}
\[
\xymatrix{1\ar[r]&\ad(G)\ar[r]&\underline{\Aut}_{S\text{-gr.}}(G)\ar[r]^p&\underline{\operatorname{Autext}}(G)\ar[r]&1.}
\]
\begin{enumeratei}
\item $\underline{\Aut}_{S\text{-gr.}}(G)$ is representable by a smooth separated $S$-scheme.
\item $\underline{\operatorname{Autext}}(G)$ is representable by a twisted constant $S$-scheme of finite generation (Exp.~X, 5.1).
\item If $G$ is splittable, the preceding exact sequence is split. More precisely, for every pinning of $G$, the morphism (cf. 1.2 ($*$))
\[
p\circ a:E_S\to\underline{\operatorname{Autext}}(G)
\]
is an isomorphism.
\end{enumeratei}
\end{theorem}

Let us first show how the theorem follows from the following lemma:
\begin{lemma}{1.4}\label{XXIV.1.4}
Under the hypotheses of (iii), $\underline{\Aut}_{S\text{-gr.}}(G)$ is the semidirect product $a(E_S)\cdot\ad(G)$.\nde{Here and in what follows we have replaced the notation $\mathrm{int}(G)$ by $\ad(G)$.}
\end{lemma}

The lemma immediately implies the theorem when $G$ is splittable. Since $G$ is locally splittable for the étale topology (Exp.~XXII, 2.3), hence also for the (fppf) topology, and since the latter is ``of effective descent'' for the fibered category of twisted constant morphisms (Exp.~X, 5.5), one deduces (ii) in the general case (cf. Exp.~IV, 4.6.8). To deduce (i), one observes that $\ad(G)$ is affine over $S$, hence the morphism $p$ is affine when $\underline{\Aut}_{S\text{-gr.}}(G)$ is representable, and one concludes by descent of affine schemes.\nde{Cf. SGA~1, VIII 2.1.}

It therefore only remains for us to prove 1.4. To do so, it suffices to prove:
\begin{lemma}{1.5}\label{XXIV.1.5}
If $(\mathscr R,\mathscr E)$ and $(\mathscr R',\mathscr E')$ are two pinnings of the reductive $S$-group $G$, there exists a unique inner automorphism $u$ of $G$ over $S$ taking one pinning to the other (i.e. such that there exists $v:\mathscr R'\isomto\mathscr R$ such that $u(\mathscr E)=v(\mathscr E')$, cf. 1.0).
\end{lemma}

\numpara{1.5.1}\label{XXIV.1.5.1}\textit{Uniqueness.} It suffices to prove that if $G$ is a pinned $S$-group and if $\mathrm{int}(g)$ is an automorphism of a pinned group ($g\in G(S)$), then $\mathrm{int}(g)=\mathrm{id}$. Now one first has $\mathrm{int}(g)T=T$, $\mathrm{int}(g)B=B$, hence $g\in\uNorm_G(T)(S)\cap\uNorm_G(B)(S)=T(S)$ (cf. for example Exp.~XXII, 5.6.1). It follows that $\mathrm{int}(g)$ normalizes each $U_\alpha$ and that
\[
\Lie(\mathrm{int}(g))X_\alpha=\Ad(g)X_\alpha=\alpha(g)X_\alpha
\]
for every $\alpha\in\Delta$. One therefore has $\alpha(g)=1$ for $\alpha\in\Delta$, hence $g\in\bigcap_{\alpha\in\Delta}(\Ker\alpha)(S)=\uCentr(G)(S)$ (Exp.~XXII, 4.1.8). \hfill Q.E.D.

% original PDF p. 5
\numpara{1.5.2}\label{XXIV.1.5.2}\textit{Existence.} It suffices to prove it locally for the (fpqc) topology. Let
\[
(T,M,R,\Delta,(X_\alpha)_{\alpha\in\Delta})\quad\text{and}\quad
(T',M',R',\Delta',(X'_{\alpha'})_{\alpha'\in\Delta'})
\]
be the two pinnings. By conjugacy of maximal tori, one may suppose $T=T'$. After restricting $S$, one may suppose that the isomorphism $\mathrm D_S(M)\simeq\mathrm D_S(M')$ comes from an isomorphism $M\simeq M'$ taking $R$ onto $R'$, and one is reduced to the situation $T=T'$, $M=M'$, $R=R'$. Since the systems of simple roots are conjugate under the Weyl group (Exp.~XXI, 3.3.7), one may also suppose $\Delta=\Delta'$. There then exists for each $\alpha\in\Delta$ a scalar $z_\alpha\in\Gm(S)$ such that $X'_\alpha=z_\alpha X_\alpha$, and it suffices to construct locally for (fpqc) a section $t$ of $T$ such that $\alpha(t)=z_\alpha$ for each $\alpha\in\Delta$. But the morphism $T\to(\mathbf G_{\mathrm m,S})^\Delta$ with components $\{\alpha,\alpha\in\Delta\}$ is the dual of an \emph{injection} $\ZZ^\Delta\to M$, hence is faithfully flat, which completes the proof of 1.5.2 and therefore of 1.3.

\begin{corollary}{1.6}\label{XXIV.1.6}
Let $S$ be a scheme and $G$ a reductive $S$-group. The following conditions are equivalent:
\begin{enumeratei}
\item $\underline{\Aut}_{S\text{-gr.}}(G)$ is affine (resp. of finite type, resp. of finite presentation, resp. quasi-compact) over $S$.
\item $\underline{\operatorname{Autext}}(G)$ is finite over $S$.
\item For every $s\in S$, one has $\operatorname{rkred}(G_s)-\operatorname{rkss}(G_s)\leq1$.
\end{enumeratei}
\end{corollary}
\begin{proof}
Indeed, since $\ad(G)$ is affine, flat and of finite presentation over $S$, the morphism $p:\underline{\Aut}_{S\text{-gr.}}(G)\to\underline{\operatorname{Autext}}(G)$ is affine, faithfully flat and of finite presentation.

If $\underline{\operatorname{Autext}}(G)$ is finite over $S$, it is affine over $S$, hence so is $\underline{\Aut}_{S\text{-gr.}}(G)$, which proves (ii) $\Rightarrow$ (i). If $\underline{\Aut}_{S\text{-gr.}}(G)$ is quasi-compact over $S$, it is of finite presentation over $S$ (being in any case locally of finite presentation and separated over $S$); by Exp.~V, 9.1, $\underline{\operatorname{Autext}}(G)$ is then of finite presentation over $S$, hence finite, which proves (i) $\Rightarrow$ (ii). Finally, to prove the equivalence of (ii) and (iii), one may suppose $G$ split, and one is reduced to Exp.~XXI, 6.7.8.
\end{proof}

\begin{corollary}{1.7}\label{XXIV.1.7}
Let $S$ be a scheme and $G$ a reductive $S$-group. Then one has $\underline{\Aut}_{S\text{-gr.}}(G)^0\simeq\ad(G)$.
\end{corollary}
\begin{corollary}{1.8}\label{XXIV.1.8}
Let $S$ be a scheme, $G$ a reductive $S$-group, and $H$ a smooth affine $S$-group scheme with connected fibers over $S$. Then the $S$-functor
\[
\underline{\Isom}_{S\text{-gr.}}(G,H)
\]
is representable by a smooth separated $S$-scheme (which is affine over $S$ if $G$ is semisimple).
\end{corollary}
Indeed, let $U$ be the set of points $s$ of $S$ such that $H_s$ is reductive; it is open (Exp.~XIX, 2.6); if $S'$ is an $S$-scheme, $H_{S'}$ is reductive if and only if $S'\to S$ factors through $U$. It follows that the canonical morphism $\underline{\Isom}_{S\text{-gr.}}(G,H)\to S$ factors through $U$. One may therefore suppose $S=U$ and one is reduced to:

% original PDF p. 6
\begin{corollary}{1.9}\label{XXIV.1.9}
Let $S$ be a scheme and $G$, $G'$ two reductive $S$-groups. Then
\[
F=\underline{\Isom}_{S\text{-gr.}}(G,G')
\]
is representable by a smooth separated $S$-scheme (affine if $G$ or $G'$ is semisimple). Moreover, $S$ decomposes as the sum of two open subschemes $S_1$ and $S_2$ such that $F_{S_1}=\varnothing$ and $F_{S_2}$ is a left (resp. right) principal homogeneous bundle under $\underline{\Aut}_{S_2\text{-gr.}}(G'_{S_2})$ (resp. $\underline{\Aut}_{S_2\text{-gr.}}(G_{S_2})$).
\end{corollary}
\begin{proof}
Indeed, let $S_2$ be the set of points of $S$ where $G$ and $G'$ have the same type, and let $S_1$ be its complement.

Since the type of a reductive group is a locally constant function, $S_1$ and $S_2$ are open. It is clear that $F_{S_1}=\varnothing$, and one may suppose $S=S_2$. By Exp.~XXIII, 5.6, $F$ is a principal homogeneous sheaf under $\underline{\Aut}_{S\text{-gr.}}(G)$, locally trivial for the étale topology. It follows that $F_0=F/\ad(G)$ is a principal homogeneous sheaf under $\underline{\operatorname{Autext}}(G)$, locally trivial for the étale topology, hence representable (Exp.~X, 5.5). Since $\ad(G)$ is affine, $F$ is therefore also representable.
\end{proof}
Let us observe that in the course of the proof one has obtained:
\begin{corollary}{1.10}\label{XXIV.1.10}
Let $S$ be a scheme and $G$, $G'$ two reductive $S$-groups of the same type at every $s\in S$. Then $\ad(G)$ acts freely (on the right) on $\underline{\Isom}_{S\text{-gr.}}(G,G')$, and the quotient sheaf
\[
\underline{\operatorname{Isomext}}(G,G')=\underline{\Isom}_{S\text{-gr.}}(G,G')/\ad(G)
\]
is representable by a twisted constant $S$-scheme, which is a principal homogeneous bundle under $\underline{\operatorname{Autext}}(G)$ (and which is therefore finite over $S$ if $G$ is semisimple). Moreover, the isomorphism
\[
\underline{\Isom}_{S\text{-gr.}}(G,G')\simeq\underline{\Isom}_{S\text{-gr.}}(G',G)
\]
defined by $u\mto u^{-1}$ induces an isomorphism
\[
\underline{\operatorname{Isomext}}(G,G')\simeq\underline{\operatorname{Isomext}}(G',G).
\]
\end{corollary}
\begin{remark}{1.11}\label{XXIV.1.11}
If $\underline{\operatorname{Isomext}}(G,G')(S)\ne\varnothing$, one says that $G$ is an inner twisted form of $G'$; then $G'$ is an inner twisted form of $G$; one can then reduce the structure group of $\underline{\Isom}_{S\text{-gr.}}(G,G')$ to $\ad(G)$. More precisely, let $u\in\underline{\operatorname{Isomext}}(G,G')(S)$, considered as a section $u:S\to\underline{\operatorname{Isomext}}(G,G')$. Denote by
\[
\underline{\operatorname{Isomint}}_u(G,G')
\]
the inverse image, under the canonical morphism $\underline{\Isom}_{S\text{-gr.}}(G,G')\to\underline{\operatorname{Isomext}}(G,G')$, of the closed subscheme of $\underline{\operatorname{Isomext}}(G,G')$ which is the image of $u$. The natural action of $\ad(G)$ on $\underline{\operatorname{Isomint}}_u(G,G')$ equips this scheme with a structure of principal homogeneous bundle; by extension of the structure group\nde{We have corrected $\ad(G)\to\underline{\operatorname{Autext}}(G)$ to $\ad(G)\to\underline{\Aut}(G)$.} $\ad(G)\to\underline{\Aut}(G)$, $\underline{\operatorname{Isomint}}_u(G,G')$ gives back $\underline{\Isom}_{S\text{-gr.}}(G,G')$.
\end{remark}
By Hensel's lemma (Exp.~XI, 1.11), 1.8 immediately gives:

% original PDF p. 7
\begin{corollary}{1.12}\label{XXIV.1.12}
Let $S$ be a henselian local scheme, $G$ a reductive $S$-group, $G'$ a smooth affine $S$-group with connected fibers, and $s$ the closed point of $S$. If $G_s$ and $G'_s$ are isomorphic algebraic $\kappa(s)$-groups, $G$ and $G'$ are isomorphic. More precisely, every $\kappa(s)$-isomorphism $G_s\simeq G'_s$ comes from an $S$-isomorphism $G\simeq G'$.
\end{corollary}
Applying 1.7 (resp. 1.12) to the scheme of dual numbers over a field, one deduces from Exp.~III, 2.10 (resp. 3.10) part (i) (resp. (ii)) of the following corollary.
\begin{corollary}{1.13}\label{XXIV.1.13}
Let $k$ be a field and $G$ a reductive $k$-group.
\begin{enumeratei}
\item If $G$ is adjoint, one has $\mathrm H^1(G,\Lie(G/k))=0$.\nde{We have corrected the original, which stated (i) without assuming $G$ adjoint. According to the characterization of the $\mathrm H^i(G,-)$ as derived functors of $\Hom_G(k,-)$ (Exp.~I, 5.3.1, see also [Ja87], I 4.16), one has $\mathrm H^i(G,V)=\Ext^i_G(k,V)$ for every $G$-module $V$; now if $\operatorname{char}(k)=p>0$ and if $G=\SL_{p,k}$, then $\Lie(\GL_p/k)$ is a nontrivial extension of $k$ by $\mathfrak g=\Lie(\SL_p/k)$, hence $\mathrm H^1(G,\mathfrak g)\ne0$. See also the addition 1.15.1 below.}
\item One has $\mathrm H^2(G,\Lie(G/k))=0$.
\end{enumeratei}
\end{corollary}
\begin{remark}{1.14}\label{XXIV.1.14}
(i) The assertion concerning $\mathrm H^1$ was known (Chevalley); that concerning $\mathrm H^2$ was proved by Chevalley in most cases of the classification.

(ii) In fact, the conjunction of 1.13 and the uniqueness theorem over an algebraically closed field is essentially equivalent to the uniqueness theorem. A direct proof of 1.13 would therefore give a way to deduce the general uniqueness theorem from Chevalley's uniqueness theorem over a field.

The existence of reductive groups of all types over all schemes (Exp.~XXV) shows that the obstructions to lifting a reductive $k$-group $G$ over Artinian rings with residue field $k$ (which, by Exp.~III, 3.8, are elements of
\[
\mathrm H^3(G,\Lie(G/k)\otimes V)\simeq\mathrm H^3(G,\Lie(G/k))\otimes V,
\]
where $V$ is a certain $k$-vector space (equipped with the trivial action of $G$)) vanish. This seems to suggest that $\mathrm H^3(G,\Lie(G/k))=0$. Here again, a direct proof of this fact (if it is true) would doubtless give a way to deduce the general existence theorem from the existence theorem over a field (Chevalley's Tôhoku).
\end{remark}
\begin{corollary}{1.15}\label{XXIV.1.15}
Let $k$ be a field,\nde{We have replaced ``scheme'' by ``field''.} and $G$ a reductive $k$-group. Consider $k$ as a trivial $G$-module. Then
\[
\mathrm H^1(G,k)=\mathrm H^2(G,k)=0.
\]
\end{corollary}
\begin{proof}
\nde{Because of the correction made in 1.13, we have given another proof in the case of $\mathrm H^1$. Moreover, it follows from a theorem of G.~Kempf that $\mathrm H^i(G,k)=0$ for every $i>0$, cf. [Ja87], II 4.5 and 4.11.}Since $\mathrm H^i(G_{\bar k},\bar k)=\mathrm H^i(G,k)\otimes\bar k$, one may suppose $k$ algebraically closed. An element of $\mathrm H^1(G,k)$ is nothing other than a morphism of $k$-groups $\phi:G\to\mathbf G_{\mathrm a,k}$. Then $\phi(G)=G/\Ker(\phi)$ is a smooth, connected and reductive subgroup (cf. XIX 1.7) of $\mathbf G_{\mathrm a,k}$, hence trivial. Thus $\mathrm H^1(G,k)=0$.

% original PDF p. 8
Consider now the reductive $k$-group $H=G\times_k\mathbf G_{\mathrm m,k}$. One has $\Lie(H/k)=\Lie(G/k)\oplus k$, a decomposition stable under $H$. For an arbitrary $H$-module $V$, one has
\[
\mathrm H^i(H,V)=\mathrm H^i(G,\mathrm H^0(\mathbf G_{\mathrm m,k},V))
\]
(this follows from the characterization of the $\mathrm H^i(H,-)$ as derived functors of $\mathrm H^0(H,-)=\mathrm H^0(G,\mathrm H^0(\mathbf G_{\mathrm m,k},-))$, and from the fact that the functor $\mathrm H^0(\mathbf G_{\mathrm m,k},-)$ is exact, cf. Exp.~I, 5.3.1 and 5.3.3). In particular, for every $i$ one has
\[
\mathrm H^i(H,\Lie(H/k))=\mathrm H^i(G,\Lie(G/k))\oplus\mathrm H^i(G,k)
\]
hence $\mathrm H^2(G,k)=0$ by applying 1.13 (ii) to the reductive group $H$.
\end{proof}
\begin{corollary}{1.15.1}\label{XXIV.1.15.1}
\nde{We have added this corollary, drawn from remarks of O.~Gabber, which makes 1.13 (i) more precise.}Let $k$ be a field, $G$ a reductive $k$-group, $Z$ its center and $\mathscr R=\allowbreak(M,\allowbreak M^*,\allowbreak R,\allowbreak R^*)$ the type of $G$. Then one has
\[
\mathrm H^1(G,\Lie(G/k))\simeq\Ext^1_{\ZZ}(M/\Gamma_0(R),k).
\]
In particular, $\mathrm H^1(G,\Lie(G/k))=0$ if and only if $Z$ is smooth over $k$.
\end{corollary}
\begin{proof}
Indeed, let $\mathfrak g$ (resp. $\mathfrak z$, resp. $\mathfrak g_{\mathrm{ad}}$) be the Lie algebra of $G$ (resp. $Z$, resp. $G_{\mathrm{ad}}=G/Z$). It follows from 1.15 (and its proof) that
\[
\mathrm H^1(G,\mathfrak g)=\mathrm H^1(G_{\mathrm{ad}},\mathfrak g)\simeq\mathrm H^1(G_{\mathrm{ad}},\mathfrak g/\mathfrak z).
\]
Put $C=\Coker(\mathfrak g\to\mathfrak g_{\mathrm{ad}})$; one has an exact sequence
\[
\xymatrix{0\ar[r]&\mathfrak g/\mathfrak z\ar[r]&\mathfrak g_{\mathrm{ad}}\ar[r]&C\ar[r]&0.}
\]
Since $\mathrm H^1(G_{\mathrm{ad}},\mathfrak g_{\mathrm{ad}})=0$ (1.13) and $\mathrm H^0(G_{\mathrm{ad}},\mathfrak g_{\mathrm{ad}})=0$ (cf. Exp.~II, 5.2.3), one obtains
\[
\mathrm H^1(G_{\mathrm{ad}},\mathfrak g/\mathfrak z)\simeq\mathrm H^0(G_{\mathrm{ad}},C).
\]
To compute the right-hand term, one may suppose $k$ algebraically closed. Let $(G,T,M,R)$ be a splitting of $G$; then $Z=\mathrm D_k(M/\Gamma_0(R))$ and $C$ identifies with
\[
\Coker\bigl(\Hom_{\ZZ}(M,k)\to\Hom_{\ZZ}(\Gamma_0(R),k)\bigr)
\simeq\Ext^1_{\ZZ}(M/\Gamma_0(R),k),
\]
equipped with the trivial action of $G_{\mathrm{ad}}$. The corollary follows, since $Z$ is not smooth over $k$ if and only if $\operatorname{char}(k)=p>0$ and $M/\Gamma_0(R)$ has $p$-torsion (Exp.~IX, 2.1).
\end{proof}
\begin{definition}{1.16}\label{XXIV.1.16}
Let $S$ be a scheme and $G$ a reductive $S$-group. One calls a \emph{form of $G$ over $S$} an $S$-group scheme $G'$ locally isomorphic to $G$ for the (fpqc) topology (it amounts to the same thing to say (cf. Exp.~XXIII, 5.6) that $G'$ is locally isomorphic to $G$ for the étale topology, or again that $G'$ is a reductive $S$-group of the same type as $G$ at every point of $S$).
\end{definition}
\begin{corollary}{1.17}\label{XXIV.1.17}
Let $S$ be a scheme and $G$ a reductive $S$-group.
\begin{enumeratei}
\item The functor
\[
G'\mto\underline{\Isom}_{S\text{-gr.}}(G,G')
\]
is an equivalence between the category of forms of $G$ over $S$ and the category of principal homogeneous bundles under $\underline{\Aut}_{S\text{-gr.}}(G)$.
% Corollary 1.17 and enumeratei continue in en-02.tex, PDF p. 9.
